%==============================================================================
% FICHAMENTO CRÍTICO — Deep Learning for Automated Detection of Osteoporosis
% on Panoramic Radiographs
% Conforme NBR 6023 (referências) e NBR 10520 (citações)
% Capítulo 8 — Osteoporose em Radiografias Panorâmicas
%==============================================================================
\section{Fichamento: \citeonline{yamamoto2024deep}}

% --- 1. IDENTIFICAÇÃO ---
\subsection{Identificação}
\begin{tabular}{ll}
\textbf{Título:} & Deep learning for automated detection of osteoporosis on
panoramic radiographs \\
\textbf{Autor(es):} & Yamamoto, Keisuke; Santos, Rafael; Park, Ji-Yeon;
Chen, Wei-Ming; Tanaka, Hiroshi; Azevedo-Marques, Paulo M. \\
\textbf{Periódico:} & Oral Surgery, Oral Medicine, Oral Pathology and Oral Radiology \\
\textbf{Ano:} & 2024 \\
\textbf{Volume:} & 138 \\
\textbf{Número:} & 4 \\
\textbf{Páginas:} & 456--470 \\
\textbf{DOI:} & \url{https://doi.org/10.1016/j.oooo.2024.01.014} \\
\textbf{Qualis:} & A2 (Odontologia) \\
\textbf{Tipo:} & Artigo original com validação externa multicêntrica \\
\end{tabular}

% --- 2. OBJETIVO ---
\subsection{Objetivo}
Desenvolver e validar um modelo de deep learning baseado em arquitetura
DenseNet-201 com módulo de atenção Squeeze-and-Excitation (SE) para
triagem automatizada de osteoporose a partir de radiografias panorâmicas,
utilizando a densitometria óssea (DXA) do colo do fêmur e coluna lombar
como padrão-ouro. O estudo visa explorar a radiografia panorâmica --
exame odontológico de baixo custo, amplamente disponível e de baixa
radiação -- como ferramenta de triagem populacional para osteoporose,
condição subdiagnosticada que afeta 200 milhões de pessoas globalmente
e cuja detecção precoce poderia prevenir fraturas osteoporóticas
graves (quadril, vértebras). O modelo analisa exclusivamente a região
da cortical mandibular inferior, onde o adelgaçamento cortical e a
perda óssea são marcadores radiomorfológicos validados de osteoporose
sistêmica (índice cortical mandibular, MCI).

% --- 3. METODOLOGIA ---
\subsection{Metodologia}
\textbf{Desenho do estudo:} Estudo retrospectivo multicêntrico de
desenvolvimento e validação de modelo de classificação binária
(osteoporose vs. não osteoporose). \\
\textbf{Amostra:} 8.942 radiografias panorâmicas digitais de pacientes
com 50 anos ou mais (74,3\% mulheres), pareadas com exame DXA realizado
em até 6 meses da panorâmica, provenientes de 7 centros clínicos em 4
países (Brasil, Japão, Coreia do Sul e Alemanha). A prevalência de
osteoporose na amostra foi de 22,8\% (DXA com T-score $\leq$ -2,5). \\
\textbf{Métricas principais:} AUC-ROC, sensibilidade, especificidade,
VPP, VPN, acurácia, F1-score, kappa de Cohen. \\
\textbf{Validação:} Validação cruzada 5-fold estratificada por centro,
sexo e faixa etária; validação externa com dataset independente
(n=1.200 panorâmicas de 2 centros não participantes do treino);
análise de subgrupos por sexo, faixa etária (50--59, 60--69, $\geq$70
anos) e estágio da osteoporose (osteopenia vs. osteoporose
estabelecida).

% --- 4. RESULTADOS PRINCIPAIS ---
\subsection{Resultados Principais}
\begin{itemize}
  \item O modelo DenseNet-201 com SE block alcançou AUC de 0,914
        (IC 95\%: 0,897--0,931) na validação interna, superando
        significativamente as baselines ResNet-50 (AUC 0,882,
        p = 0,003) e EfficientNet-B4 (AUC 0,871, p < 0,001), bem
        como o índice cortical mandibular manual (MCI kappa 0,68,
        AUC 0,793).
  \item Sensibilidade de 0,887 e especificidade de 0,903 no threshold
        de Youden, com VPP de 0,734 e VPN de 0,961, indicando que
        o modelo é excelente para descartar osteoporose (VPN > 0,95)
        e moderado para confirmar a doença (VPP limitado pela
        prevalência de 22,8\%).
  \item Na validação externa independente (n=1.200), a AUC foi de
        0,892 (IC 95\%: 0,871--0,913), 2,2 pontos percentuais abaixo
        da validação interna, confirmando boa generalização.
  \item A análise de subgrupos revelou AUC de 0,928 para mulheres
        acima de 60 anos (subgrupo de maior risco), mas queda
        significativa para homens (AUC 0,842), sugerindo que a
        osteoporose masculina, menos estudada e com manifestações
        radiográficas distintas, representa um desafio adicional.
  \item A sensibilidade para osteopenia (T-score entre -1,0 e -2,5)
        foi de 0,721, substancialmente inferior à sensibilidade para
        osteoporose estabelecida (0,887), indicando que o modelo é
        menos sensível para estágios iniciais da perda óssea.
  \item O Grad-CAM++ aplicado às camadas de atenção SE revelou que as
        regiões mais ativadas correspondem consistentemente à borda
        inferior da mandíbula (cortical mandibular), córtex do
        côndilo e assoalho do seio maxilar, confirmando que o modelo
        aprendeu características radiomorfológicas alinhadas ao
        conhecimento radiológico estabelecido.
\end{itemize}

% --- 5. LIMITAÇÕES ---
\subsection{Limitações}
\begin{itemize}
  \item A sensibilidade reduzida para osteopenia (0,721) limita a
        utilidade do modelo como ferramenta de triagem populacional
        para estágios iniciais, quando intervenções preventivas
        (suplementação, exercício, medicação) são mais eficazes.
  \item O desempenho inferior em homens (AUC 0,842 vs. 0,928 em
        mulheres) é clinicamente relevante, pois fraturas
        osteoporóticas em homens têm maior morbimortalidade e a
        osteoporose masculina é particularmente subdiagnosticada.
  \item A qualidade e o padrão da radiografia panorâmica variam
        significativamente entre equipamentos e técnicos; o modelo
        foi treinado majoritariamente com imagens de equipamentos
        digitais de fabricantes específicos (Sirona, Planmeca), e
        imagens de equipamentos mais antigos ou analógicos
        digitalizados não foram testadas.
  \item O pareamento com DXA foi realizado com janela de até 6 meses,
        durante a qual o status ósseo do paciente pode ter se
        alterado, especialmente em pacientes em tratamento para
        osteoporose, introduzindo erro de rótulo (label noise).
  \item O estudo não incluiu pacientes com fraturas osteoporóticas
        prévias, artropatias degenerativas avançadas da ATM ou
        osteomielite mandibular, condições que alteram a morfologia
        óssea e poderiam gerar falso-positivos.
\end{itemize}

% --- 6. CONTRIBUIÇÃO PARA O LIVRO ---
\subsection{Contribuição para o Livro}
Este artigo fundamenta o Capítulo 8, que aborda a aplicação de IA para
triagem de osteoporose em radiografias panorâmicas. O modelo DenseNet-201
com SE block é utilizado no livro como estudo de caso de mecanismos de
atenção aplicados à imagem odontológica, demonstrando como a atenção
Squeeze-and-Excitation pode melhorar a representação de características
discriminativas em exames de baixo contraste. A comparação com o índice
cortical mandibular manual (MCI) ilustra o argumento do livro sobre a
superioridade de modelos de aprendizado de características
(feature learning) sobre índices radiomorfológicos tradicionais. A
distinção entre sensibilidade para osteoporose estabelecida vs.
osteopenia é discutida no livro como exemplo de maturidade tecnológica
(TRL) -- o modelo é maduro para triagem de casos avançados mas requer
melhorias para detecção precoce. Os mapas Grad-CAM++ são empregados
como exemplo de validação clínica de modelos de IA.

% --- 7. ANÁLISE CRÍTICA ---
\subsection{Análise Crítica}
\begin{itemize}
  \item \textbf{Validade interna:} Alta -- Validação cruzada estratificada
        5-fold com repetição; comparação sistemática com múltiplas
        arquiteturas de CNN; threshold de Youden para otimização
        clínica; análise de calibração.
  \item \textbf{Validade externa:} Alta -- Dataset independente de 1.200
        imagens de 2 centros não participantes do treino, com
        diferentes equipamentos e distribuição demográfica distinta
        (população predominantemente caucasiana na validação externa
        vs. asiática no treino).
  \item \textbf{Reprodutibilidade:} Média-Baixa -- Embora o pipeline de
        pré-processamento e a arquitetura DenseNet-201 sejam públicos,
        os pesos do modelo treinado e o dataset anotado não foram
        disponibilizados; a dependência de DXA como padrão-ouro
        limita a reprodutibilidade em contextos sem acesso ao exame.
  \item \textbf{Relevância clínica:} Muito Alta -- A radiografia
        panorâmica é o exame odontológico mais realizado globalmente;
        utilizar uma imagem já disponível para triagem de osteoporose
        sem custo adicional de radiação ou exame representa um ganho
        populacional imenso, especialmente em países em
        desenvolvimento com acesso limitado a DXA.
  \item \textbf{Conflito de interesses:} Não -- Todos os autores
        declararam ausência de conflitos de interesse; o estudo foi
        financiado por agências de fomento públicas (CNPq, JSPS,
        DFG) sem participação da indústria.
  \item \textbf{Viés identificado:} Viés de confirmação -- a amostra
        incluiu apenas pacientes com DXA disponível, o que pode
        super-representar pacientes com suspeita clínica de
        osteoporose (viés de encaminhamento), com prevalência de
        22,8\% superior aos 8--12\% esperados na população geral
        acima de 50 anos, inflando o VPP; viés de rótulo -- janela
        de 6 meses entre panorâmica e DXA pode gerar erro de
        classificação em pacientes em tratamento.
\end{itemize}

% --- 8. CITAÇÕES RELEVANTES ---
\subsection{Citações Relevantes}
\begin{citacao}
``The DenseNet-201 model with Squeeze-and-Excitation attention blocks
achieved an AUC of 0.914 (95\% CI 0.897--0.931) for osteoporosis
detection on panoramic radiographs, significantly outperforming the
manual mandibular cortical index (MCI, AUC 0.793). Sensitivity was 0.887
for established osteoporosis but dropped to 0.721 for osteopenia,
indicating stage-dependent performance. External validation on 1,200
independent images yielded an AUC of 0.892 (95\% CI 0.871--0.913).''
(p. 464).
\end{citacao}

% --- 9. MAPEAMENTO SDD ---
\subsection{Mapeamento SDD}
\textbf{Problema:} Osteoporose é subdiagnosticada globalmente (estima-se
que 80\% dos casos não são detectados), e a DXA -- padrão-ouro -- tem
acesso limitado em países em desenvolvimento. Radiografias panorâmicas
são amplamente disponíveis na prática odontológica, mas a avaliação
visual da cortical mandibular tem acurácia limitada (AUC $\sim$0,79)
e variabilidade interexaminar relevante. \\
\textbf{Entrada:} Radiografias panorâmicas digitais (8.942 imagens, 7
centros, 4 países); pré-processamento com equalização de histograma,
normalização e recorte da região mandibular; pareamento com DXA do
fêmur/coluna (T-score). \\
\textbf{Saída:} Probabilidade de osteoporose (T-score $\leq$ -2,5) com
classificação binária; mapa de ativação Grad-CAM++ sobreposto à
panorâmica para interpretabilidade clínica. \\
\textbf{Critérios:} AUC $\geq$ 0,90; sensibilidade $\geq$ 0,85 para
osteoporose estabelecida; VPN $\geq$ 0,95 (priorizando não deixar
casos positivos sem encaminhamento); degradação em validação externa
$\leq$ 5\%; tempo de inferência $\leq$ 5s/imagem. \\
\textbf{Confiança:} 0,84 -- Estudo multicêntrico robusto com validação
externa independente, padrão-ouro DXA e mapas de ativação consistentes
com a anatomia, mas limitado pela sensibilidade reduzida para osteopenia
e desempenho inferior em homens, que restringem o escopo de aplicação
para triagem populacional universal.

% --- 10. REFERÊNCIA ABNT ---
\subsection{Referência ABNT}
\begin{verbatim}
YAMAMOTO, Keisuke et al. Deep learning for automated detection of
osteoporosis on panoramic radiographs. Oral Surgery, Oral Medicine, Oral
Pathology and Oral Radiology, New York, v. 138, n. 4, p. 456-470,
2024. DOI: 10.1016/j.oooo.2024.01.014.
\end{verbatim}
\endinput
